% paper58-refactoring-guidance.tex — Sheaf-Theoretic Refactoring: When Can Code Be Safely Restructured?
%   Paper 58 of the Judgment Geometry series.
% Compile: pdflatex paper58-refactoring-guidance
\documentclass[11pt]{article}
\usepackage{lmodern}
% jugeo-common.tex — shared preamble for all Judgment Geometry papers
% Include via: % jugeo-common.tex — shared preamble for all Judgment Geometry papers
% Include via: % jugeo-common.tex — shared preamble for all Judgment Geometry papers
% Include via: \input{jugeo-common}

% ── Packages ────────────────────────────────────────────────────────────
\usepackage{amsmath,amssymb,amsthm,mathtools}
\usepackage{tikz-cd}
\usepackage{listings}
\usepackage{xcolor}
\usepackage{xspace}
\usepackage{booktabs}
\usepackage{multirow}
\usepackage{enumitem}
\usepackage[expansion=false]{microtype}
\usepackage{hyperref}
\usepackage{cleveref}
% \usepackage{thmtools}  % disabled: not available in this TeX installation
\usepackage{float}
\usepackage{caption}
\usepackage{subcaption}
\usepackage{tabularx}
\usepackage{stmaryrd}       % \llbracket, \rrbracket
\usepackage{bbm}            % \mathbbm{1}

% ── Page geometry ───────────────────────────────────────────────────────
\usepackage[margin=1in]{geometry}

% ── Theorem environments ────────────────────────────────────────────────
\theoremstyle{plain}
\newtheorem{theorem}{Theorem}[section]
\newtheorem{lemma}[theorem]{Lemma}
\newtheorem{proposition}[theorem]{Proposition}
\newtheorem{corollary}[theorem]{Corollary}

\theoremstyle{definition}
\newtheorem{definition}[theorem]{Definition}
\newtheorem{example}[theorem]{Example}
\newtheorem{construction}[theorem]{Construction}

\theoremstyle{remark}
\newtheorem{remark}[theorem]{Remark}
\newtheorem{notation}[theorem]{Notation}

% ── Python listings style ──────────────────────────────────────────────
\definecolor{jugeoBlue}{HTML}{2563EB}
\definecolor{jugeoGreen}{HTML}{059669}
\definecolor{jugeoRed}{HTML}{DC2626}
\definecolor{jugeoPurple}{HTML}{7C3AED}
\definecolor{jugeoGray}{HTML}{6B7280}
\definecolor{jugeoBg}{HTML}{F8FAFC}

\lstdefinestyle{jugeo-python}{
  language=Python,
  basicstyle=\ttfamily\small,
  keywordstyle=\color{jugeoBlue}\bfseries,
  stringstyle=\color{jugeoGreen},
  commentstyle=\color{jugeoGray}\itshape,
  numberstyle=\tiny\color{jugeoGray},
  backgroundcolor=\color{jugeoBg},
  numbers=left,
  numbersep=6pt,
  frame=single,
  framerule=0.4pt,
  rulecolor=\color{jugeoGray!40},
  breaklines=true,
  breakatwhitespace=true,
  showstringspaces=false,
  tabsize=4,
  xleftmargin=1.5em,
  framexleftmargin=1.5em,
  morekeywords={dataclass,frozen,field,Enum,IntEnum,auto,Any,
                Mapping,Sequence,Optional,match,case},
  literate={→}{{$\to$}}1 {←}{{$\leftarrow$}}1
           {≤}{{$\leq$}}1 {≥}{{$\geq$}}1 {≠}{{$\neq$}}1
           {∈}{{$\in$}}1 {∀}{{$\forall$}}1 {∃}{{$\exists$}}1
           {⊕}{{$\oplus$}}1 {⊖}{{$\ominus$}}1,
  escapeinside={(*@}{@*)},
}
\lstset{style=jugeo-python}

\lstdefinestyle{jugeo-output}{
  basicstyle=\ttfamily\small,
  backgroundcolor=\color{jugeoBg},
  frame=single,
  framerule=0.4pt,
  rulecolor=\color{jugeoGray!40},
  breaklines=true,
  showstringspaces=false,
  xleftmargin=1.5em,
  framexleftmargin=0.5em,
  numbers=none,
}

% ── JuGeo-specific macros ──────────────────────────────────────────────

% Judgment 8-tuple
\newcommand{\J}{\mathcal{J}}
\newcommand{\Jud}[8]{(#1,\,#2,\,#3,\,#4,\,#5,\,#6,\,#7,\,#8)}
\newcommand{\JudShort}[1]{\J_{#1}}

% Components
\newcommand{\coord}{\mathsf{c}}            % coordinate
\newcommand{\prop}{\varphi}                 % proposition
\newcommand{\carrier}{\mathcal{A}}          % carrier type
\newcommand{\evid}{\mathcal{E}}             % evidence bundle
\newcommand{\oblig}{\mathcal{O}}            % obligations
\newcommand{\obstr}{\mathcal{B}}            % obstructions
\newcommand{\trust}{\mathcal{T}}            % trust annotation
\newcommand{\prov}{\Pi}                     % provenance

% Trust algebra
\newcommand{\Talg}{\mathfrak{T}}            % trust ordered algebra
\newcommand{\Eadm}{\mathcal{E}_{\mathrm{adm}}}
\newcommand{\tleq}{\preceq}                 % trust ordering
\newcommand{\tjoin}{\oplus}                 % conservative join
\newcommand{\tatten}{\ominus}               % attenuation
\newcommand{\tpromote}{\uparrow_\pi}        % promotion
\newcommand{\tdemote}{\downarrow_\chi}       % demotion

% Trust tiers
\newcommand{\Tcontra}{\ensuremath{\mathsf{CONTRADICTED}}}
\newcommand{\Tunver}{\ensuremath{\mathsf{UNVERIFIED}}}
\newcommand{\Tcopilot}{\ensuremath{\mathsf{COPILOT}}}
\newcommand{\Toracle}{\ensuremath{\mathsf{ORACLE}}}
\newcommand{\Truntime}{\ensuremath{\mathsf{RUNTIME}}}
\newcommand{\Tsolver}{\ensuremath{\mathsf{SOLVER}}}
\newcommand{\Tproof}{\ensuremath{\mathsf{PROOF}}}

% Site and geometry
\newcommand{\Site}{\mathcal{S}}             % semantic site
\newcommand{\Cov}{\mathrm{Cov}}            % covering family
\newcommand{\Ob}{\mathrm{Ob}}              % objects
\newcommand{\Mor}{\mathrm{Mor}}            % morphisms
\newcommand{\res}{\mathrm{res}}            % restriction
\newcommand{\Cech}{\check{C}}              % Čech complex
\newcommand{\Hcoh}[1]{H^{#1}}             % cohomology group

% Descent
\newcommand{\descent}{\mathrm{desc}}
\newcommand{\glue}{\mathrm{glue}}
\newcommand{\overlap}[2]{#1 \cap #2}

% Semantic moves
\newcommand{\VERIFY}{\textsc{Verify}}
\newcommand{\CONSTRUCT}{\textsc{Construct}}
\newcommand{\NORMALIZE}{\textsc{Normalize}}
\newcommand{\GLUE}{\textsc{Glue}}
\newcommand{\RESTRICT}{\textsc{Restrict}}
\newcommand{\TRANSPORT}{\textsc{Transport}}
\newcommand{\REPAIR}{\textsc{Repair}}
\newcommand{\PROMOTE}{\textsc{Promote}}
\newcommand{\CHALLENGE}{\textsc{Challenge}}
\newcommand{\ESCALATE}{\textsc{Escalate}}

% Fragments
\newcommand{\QFLIA}{\texttt{QF\_LIA}}
\newcommand{\QFLRA}{\texttt{QF\_LRA}}
\newcommand{\QFBV}{\texttt{QF\_BV}}
\newcommand{\QFUF}{\texttt{QF\_UF}}

% Misc
\newcommand{\Python}{\textsc{Python}}
\newcommand{\JuGeo}{\textsc{JuGeo}}
\newcommand{\LEAN}{\textsc{Lean}\,4}
\newcommand{\Fstar}{F$^{\star}$}
\newcommand{\Coq}{\textsc{Coq}}
\newcommand{\Dafny}{\textsc{Dafny}}

% Common shortcuts
\newcommand{\ie}{i.e.\@\xspace}
\newcommand{\eg}{e.g.\@\xspace}
\newcommand{\cf}{cf.\@\xspace}
\newcommand{\wrt}{w.r.t.\@\xspace}

% Evaluation shortcuts
\newcommand{\numcases}{300}
\newcommand{\numspec}{100}
\newcommand{\numeq}{100}
\newcommand{\numbug}{100}
\newcommand{\medtime}{6.5\,ms}
\newcommand{\totaltime}{1.949\,s}


% ── Packages ────────────────────────────────────────────────────────────
\usepackage{amsmath,amssymb,amsthm,mathtools}
\usepackage{tikz-cd}
\usepackage{listings}
\usepackage{xcolor}
\usepackage{xspace}
\usepackage{booktabs}
\usepackage{multirow}
\usepackage{enumitem}
\usepackage[expansion=false]{microtype}
\usepackage{hyperref}
\usepackage{cleveref}
% \usepackage{thmtools}  % disabled: not available in this TeX installation
\usepackage{float}
\usepackage{caption}
\usepackage{subcaption}
\usepackage{tabularx}
\usepackage{stmaryrd}       % \llbracket, \rrbracket
\usepackage{bbm}            % \mathbbm{1}

% ── Page geometry ───────────────────────────────────────────────────────
\usepackage[margin=1in]{geometry}

% ── Theorem environments ────────────────────────────────────────────────
\theoremstyle{plain}
\newtheorem{theorem}{Theorem}[section]
\newtheorem{lemma}[theorem]{Lemma}
\newtheorem{proposition}[theorem]{Proposition}
\newtheorem{corollary}[theorem]{Corollary}

\theoremstyle{definition}
\newtheorem{definition}[theorem]{Definition}
\newtheorem{example}[theorem]{Example}
\newtheorem{construction}[theorem]{Construction}

\theoremstyle{remark}
\newtheorem{remark}[theorem]{Remark}
\newtheorem{notation}[theorem]{Notation}

% ── Python listings style ──────────────────────────────────────────────
\definecolor{jugeoBlue}{HTML}{2563EB}
\definecolor{jugeoGreen}{HTML}{059669}
\definecolor{jugeoRed}{HTML}{DC2626}
\definecolor{jugeoPurple}{HTML}{7C3AED}
\definecolor{jugeoGray}{HTML}{6B7280}
\definecolor{jugeoBg}{HTML}{F8FAFC}

\lstdefinestyle{jugeo-python}{
  language=Python,
  basicstyle=\ttfamily\small,
  keywordstyle=\color{jugeoBlue}\bfseries,
  stringstyle=\color{jugeoGreen},
  commentstyle=\color{jugeoGray}\itshape,
  numberstyle=\tiny\color{jugeoGray},
  backgroundcolor=\color{jugeoBg},
  numbers=left,
  numbersep=6pt,
  frame=single,
  framerule=0.4pt,
  rulecolor=\color{jugeoGray!40},
  breaklines=true,
  breakatwhitespace=true,
  showstringspaces=false,
  tabsize=4,
  xleftmargin=1.5em,
  framexleftmargin=1.5em,
  morekeywords={dataclass,frozen,field,Enum,IntEnum,auto,Any,
                Mapping,Sequence,Optional,match,case},
  literate={→}{{$\to$}}1 {←}{{$\leftarrow$}}1
           {≤}{{$\leq$}}1 {≥}{{$\geq$}}1 {≠}{{$\neq$}}1
           {∈}{{$\in$}}1 {∀}{{$\forall$}}1 {∃}{{$\exists$}}1
           {⊕}{{$\oplus$}}1 {⊖}{{$\ominus$}}1,
  escapeinside={(*@}{@*)},
}
\lstset{style=jugeo-python}

\lstdefinestyle{jugeo-output}{
  basicstyle=\ttfamily\small,
  backgroundcolor=\color{jugeoBg},
  frame=single,
  framerule=0.4pt,
  rulecolor=\color{jugeoGray!40},
  breaklines=true,
  showstringspaces=false,
  xleftmargin=1.5em,
  framexleftmargin=0.5em,
  numbers=none,
}

% ── JuGeo-specific macros ──────────────────────────────────────────────

% Judgment 8-tuple
\newcommand{\J}{\mathcal{J}}
\newcommand{\Jud}[8]{(#1,\,#2,\,#3,\,#4,\,#5,\,#6,\,#7,\,#8)}
\newcommand{\JudShort}[1]{\J_{#1}}

% Components
\newcommand{\coord}{\mathsf{c}}            % coordinate
\newcommand{\prop}{\varphi}                 % proposition
\newcommand{\carrier}{\mathcal{A}}          % carrier type
\newcommand{\evid}{\mathcal{E}}             % evidence bundle
\newcommand{\oblig}{\mathcal{O}}            % obligations
\newcommand{\obstr}{\mathcal{B}}            % obstructions
\newcommand{\trust}{\mathcal{T}}            % trust annotation
\newcommand{\prov}{\Pi}                     % provenance

% Trust algebra
\newcommand{\Talg}{\mathfrak{T}}            % trust ordered algebra
\newcommand{\Eadm}{\mathcal{E}_{\mathrm{adm}}}
\newcommand{\tleq}{\preceq}                 % trust ordering
\newcommand{\tjoin}{\oplus}                 % conservative join
\newcommand{\tatten}{\ominus}               % attenuation
\newcommand{\tpromote}{\uparrow_\pi}        % promotion
\newcommand{\tdemote}{\downarrow_\chi}       % demotion

% Trust tiers
\newcommand{\Tcontra}{\ensuremath{\mathsf{CONTRADICTED}}}
\newcommand{\Tunver}{\ensuremath{\mathsf{UNVERIFIED}}}
\newcommand{\Tcopilot}{\ensuremath{\mathsf{COPILOT}}}
\newcommand{\Toracle}{\ensuremath{\mathsf{ORACLE}}}
\newcommand{\Truntime}{\ensuremath{\mathsf{RUNTIME}}}
\newcommand{\Tsolver}{\ensuremath{\mathsf{SOLVER}}}
\newcommand{\Tproof}{\ensuremath{\mathsf{PROOF}}}

% Site and geometry
\newcommand{\Site}{\mathcal{S}}             % semantic site
\newcommand{\Cov}{\mathrm{Cov}}            % covering family
\newcommand{\Ob}{\mathrm{Ob}}              % objects
\newcommand{\Mor}{\mathrm{Mor}}            % morphisms
\newcommand{\res}{\mathrm{res}}            % restriction
\newcommand{\Cech}{\check{C}}              % Čech complex
\newcommand{\Hcoh}[1]{H^{#1}}             % cohomology group

% Descent
\newcommand{\descent}{\mathrm{desc}}
\newcommand{\glue}{\mathrm{glue}}
\newcommand{\overlap}[2]{#1 \cap #2}

% Semantic moves
\newcommand{\VERIFY}{\textsc{Verify}}
\newcommand{\CONSTRUCT}{\textsc{Construct}}
\newcommand{\NORMALIZE}{\textsc{Normalize}}
\newcommand{\GLUE}{\textsc{Glue}}
\newcommand{\RESTRICT}{\textsc{Restrict}}
\newcommand{\TRANSPORT}{\textsc{Transport}}
\newcommand{\REPAIR}{\textsc{Repair}}
\newcommand{\PROMOTE}{\textsc{Promote}}
\newcommand{\CHALLENGE}{\textsc{Challenge}}
\newcommand{\ESCALATE}{\textsc{Escalate}}

% Fragments
\newcommand{\QFLIA}{\texttt{QF\_LIA}}
\newcommand{\QFLRA}{\texttt{QF\_LRA}}
\newcommand{\QFBV}{\texttt{QF\_BV}}
\newcommand{\QFUF}{\texttt{QF\_UF}}

% Misc
\newcommand{\Python}{\textsc{Python}}
\newcommand{\JuGeo}{\textsc{JuGeo}}
\newcommand{\LEAN}{\textsc{Lean}\,4}
\newcommand{\Fstar}{F$^{\star}$}
\newcommand{\Coq}{\textsc{Coq}}
\newcommand{\Dafny}{\textsc{Dafny}}

% Common shortcuts
\newcommand{\ie}{i.e.\@\xspace}
\newcommand{\eg}{e.g.\@\xspace}
\newcommand{\cf}{cf.\@\xspace}
\newcommand{\wrt}{w.r.t.\@\xspace}

% Evaluation shortcuts
\newcommand{\numcases}{300}
\newcommand{\numspec}{100}
\newcommand{\numeq}{100}
\newcommand{\numbug}{100}
\newcommand{\medtime}{6.5\,ms}
\newcommand{\totaltime}{1.949\,s}


% ── Packages ────────────────────────────────────────────────────────────
\usepackage{amsmath,amssymb,amsthm,mathtools}
\usepackage{tikz-cd}
\usepackage{listings}
\usepackage{xcolor}
\usepackage{xspace}
\usepackage{booktabs}
\usepackage{multirow}
\usepackage{enumitem}
\usepackage[expansion=false]{microtype}
\usepackage{hyperref}
\usepackage{cleveref}
% \usepackage{thmtools}  % disabled: not available in this TeX installation
\usepackage{float}
\usepackage{caption}
\usepackage{subcaption}
\usepackage{tabularx}
\usepackage{stmaryrd}       % \llbracket, \rrbracket
\usepackage{bbm}            % \mathbbm{1}

% ── Page geometry ───────────────────────────────────────────────────────
\usepackage[margin=1in]{geometry}

% ── Theorem environments ────────────────────────────────────────────────
\theoremstyle{plain}
\newtheorem{theorem}{Theorem}[section]
\newtheorem{lemma}[theorem]{Lemma}
\newtheorem{proposition}[theorem]{Proposition}
\newtheorem{corollary}[theorem]{Corollary}

\theoremstyle{definition}
\newtheorem{definition}[theorem]{Definition}
\newtheorem{example}[theorem]{Example}
\newtheorem{construction}[theorem]{Construction}

\theoremstyle{remark}
\newtheorem{remark}[theorem]{Remark}
\newtheorem{notation}[theorem]{Notation}

% ── Python listings style ──────────────────────────────────────────────
\definecolor{jugeoBlue}{HTML}{2563EB}
\definecolor{jugeoGreen}{HTML}{059669}
\definecolor{jugeoRed}{HTML}{DC2626}
\definecolor{jugeoPurple}{HTML}{7C3AED}
\definecolor{jugeoGray}{HTML}{6B7280}
\definecolor{jugeoBg}{HTML}{F8FAFC}

\lstdefinestyle{jugeo-python}{
  language=Python,
  basicstyle=\ttfamily\small,
  keywordstyle=\color{jugeoBlue}\bfseries,
  stringstyle=\color{jugeoGreen},
  commentstyle=\color{jugeoGray}\itshape,
  numberstyle=\tiny\color{jugeoGray},
  backgroundcolor=\color{jugeoBg},
  numbers=left,
  numbersep=6pt,
  frame=single,
  framerule=0.4pt,
  rulecolor=\color{jugeoGray!40},
  breaklines=true,
  breakatwhitespace=true,
  showstringspaces=false,
  tabsize=4,
  xleftmargin=1.5em,
  framexleftmargin=1.5em,
  morekeywords={dataclass,frozen,field,Enum,IntEnum,auto,Any,
                Mapping,Sequence,Optional,match,case},
  literate={→}{{$\to$}}1 {←}{{$\leftarrow$}}1
           {≤}{{$\leq$}}1 {≥}{{$\geq$}}1 {≠}{{$\neq$}}1
           {∈}{{$\in$}}1 {∀}{{$\forall$}}1 {∃}{{$\exists$}}1
           {⊕}{{$\oplus$}}1 {⊖}{{$\ominus$}}1,
  escapeinside={(*@}{@*)},
}
\lstset{style=jugeo-python}

\lstdefinestyle{jugeo-output}{
  basicstyle=\ttfamily\small,
  backgroundcolor=\color{jugeoBg},
  frame=single,
  framerule=0.4pt,
  rulecolor=\color{jugeoGray!40},
  breaklines=true,
  showstringspaces=false,
  xleftmargin=1.5em,
  framexleftmargin=0.5em,
  numbers=none,
}

% ── JuGeo-specific macros ──────────────────────────────────────────────

% Judgment 8-tuple
\newcommand{\J}{\mathcal{J}}
\newcommand{\Jud}[8]{(#1,\,#2,\,#3,\,#4,\,#5,\,#6,\,#7,\,#8)}
\newcommand{\JudShort}[1]{\J_{#1}}

% Components
\newcommand{\coord}{\mathsf{c}}            % coordinate
\newcommand{\prop}{\varphi}                 % proposition
\newcommand{\carrier}{\mathcal{A}}          % carrier type
\newcommand{\evid}{\mathcal{E}}             % evidence bundle
\newcommand{\oblig}{\mathcal{O}}            % obligations
\newcommand{\obstr}{\mathcal{B}}            % obstructions
\newcommand{\trust}{\mathcal{T}}            % trust annotation
\newcommand{\prov}{\Pi}                     % provenance

% Trust algebra
\newcommand{\Talg}{\mathfrak{T}}            % trust ordered algebra
\newcommand{\Eadm}{\mathcal{E}_{\mathrm{adm}}}
\newcommand{\tleq}{\preceq}                 % trust ordering
\newcommand{\tjoin}{\oplus}                 % conservative join
\newcommand{\tatten}{\ominus}               % attenuation
\newcommand{\tpromote}{\uparrow_\pi}        % promotion
\newcommand{\tdemote}{\downarrow_\chi}       % demotion

% Trust tiers
\newcommand{\Tcontra}{\ensuremath{\mathsf{CONTRADICTED}}}
\newcommand{\Tunver}{\ensuremath{\mathsf{UNVERIFIED}}}
\newcommand{\Tcopilot}{\ensuremath{\mathsf{COPILOT}}}
\newcommand{\Toracle}{\ensuremath{\mathsf{ORACLE}}}
\newcommand{\Truntime}{\ensuremath{\mathsf{RUNTIME}}}
\newcommand{\Tsolver}{\ensuremath{\mathsf{SOLVER}}}
\newcommand{\Tproof}{\ensuremath{\mathsf{PROOF}}}

% Site and geometry
\newcommand{\Site}{\mathcal{S}}             % semantic site
\newcommand{\Cov}{\mathrm{Cov}}            % covering family
\newcommand{\Ob}{\mathrm{Ob}}              % objects
\newcommand{\Mor}{\mathrm{Mor}}            % morphisms
\newcommand{\res}{\mathrm{res}}            % restriction
\newcommand{\Cech}{\check{C}}              % Čech complex
\newcommand{\Hcoh}[1]{H^{#1}}             % cohomology group

% Descent
\newcommand{\descent}{\mathrm{desc}}
\newcommand{\glue}{\mathrm{glue}}
\newcommand{\overlap}[2]{#1 \cap #2}

% Semantic moves
\newcommand{\VERIFY}{\textsc{Verify}}
\newcommand{\CONSTRUCT}{\textsc{Construct}}
\newcommand{\NORMALIZE}{\textsc{Normalize}}
\newcommand{\GLUE}{\textsc{Glue}}
\newcommand{\RESTRICT}{\textsc{Restrict}}
\newcommand{\TRANSPORT}{\textsc{Transport}}
\newcommand{\REPAIR}{\textsc{Repair}}
\newcommand{\PROMOTE}{\textsc{Promote}}
\newcommand{\CHALLENGE}{\textsc{Challenge}}
\newcommand{\ESCALATE}{\textsc{Escalate}}

% Fragments
\newcommand{\QFLIA}{\texttt{QF\_LIA}}
\newcommand{\QFLRA}{\texttt{QF\_LRA}}
\newcommand{\QFBV}{\texttt{QF\_BV}}
\newcommand{\QFUF}{\texttt{QF\_UF}}

% Misc
\newcommand{\Python}{\textsc{Python}}
\newcommand{\JuGeo}{\textsc{JuGeo}}
\newcommand{\LEAN}{\textsc{Lean}\,4}
\newcommand{\Fstar}{F$^{\star}$}
\newcommand{\Coq}{\textsc{Coq}}
\newcommand{\Dafny}{\textsc{Dafny}}

% Common shortcuts
\newcommand{\ie}{i.e.\@\xspace}
\newcommand{\eg}{e.g.\@\xspace}
\newcommand{\cf}{cf.\@\xspace}
\newcommand{\wrt}{w.r.t.\@\xspace}

% Evaluation shortcuts
\newcommand{\numcases}{300}
\newcommand{\numspec}{100}
\newcommand{\numeq}{100}
\newcommand{\numbug}{100}
\newcommand{\medtime}{6.5\,ms}
\newcommand{\totaltime}{1.949\,s}

% experiment-data.tex — AUTO-GENERATED from real jugeo CLI runs
% DO NOT EDIT — regenerate with: python3 experiments/run_universal_metrics.py
% Generated from 30 programs across 6 domains

\newcommand{\expTotalPrograms}{30}
\newcommand{\expVerified}{30}
\newcommand{\expAccuracy}{100.0\%}
\newcommand{\expCoordsMin}{2}
\newcommand{\expCoordsMax}{10}
\newcommand{\expCoordsMean}{5.2}
\newcommand{\expCoordsMedian}{5.5}
\newcommand{\expPropsMin}{4}
\newcommand{\expPropsMax}{20}
\newcommand{\expPropsMean}{10.4}
\newcommand{\expPropsSum}{311}
\newcommand{\expPropsOkSum}{310}
\newcommand{\expTimeMin}{3.506\,s}
\newcommand{\expTimeMax}{6.714\,s}
\newcommand{\expTimeMean}{4.64\,s}
\newcommand{\expTimeMedian}{4.508\,s}
\newcommand{\expTimeTotal}{139.204\,s}
\newcommand{\expObstructionTotal}{0}
\newcommand{\expHOneZero}{30}
\newcommand{\expDomainCount}{6}
\newcommand{\expTrustCOPILOTSUGGESTED}{30}
\newcommand{\expDomainDsCount}{5}
\newcommand{\expDomainDsVerified}{5}
\newcommand{\expDomainDsAvgCoords}{7.6}
\newcommand{\expDomainDsAvgProps}{15.2}
\newcommand{\expDomainMathCount}{5}
\newcommand{\expDomainMathVerified}{5}
\newcommand{\expDomainMathAvgCoords}{4.8}
\newcommand{\expDomainMathAvgProps}{9.4}
\newcommand{\expDomainSmCount}{5}
\newcommand{\expDomainSmVerified}{5}
\newcommand{\expDomainSmAvgCoords}{7.6}
\newcommand{\expDomainSmAvgProps}{15.2}
\newcommand{\expDomainSortCount}{5}
\newcommand{\expDomainSortVerified}{5}
\newcommand{\expDomainSortAvgCoords}{2.4}
\newcommand{\expDomainSortAvgProps}{4.8}
\newcommand{\expDomainStrCount}{5}
\newcommand{\expDomainStrVerified}{5}
\newcommand{\expDomainStrAvgCoords}{3.2}
\newcommand{\expDomainStrAvgProps}{6.4}
\newcommand{\expDomainWebCount}{5}
\newcommand{\expDomainWebVerified}{5}
\newcommand{\expDomainWebAvgCoords}{5.6}
\newcommand{\expDomainWebAvgProps}{11.2}

% subsystem-data.tex — AUTO-GENERATED from real jugeo subsystem experiments
% DO NOT EDIT — regenerate with: python3 experiments/run_subsystem_experiments.py

\newcommand{\subCliTotal}{10}
\newcommand{\subCliVerified}{9}
\newcommand{\subCliAccuracy}{90.0\%}
\newcommand{\subCliMeanTime}{4.132\,s}
\newcommand{\subCliMinTime}{3.472\,s}
\newcommand{\subCliMaxTime}{5.372\,s}
\newcommand{\subCliTotalCoords}{54}
\newcommand{\subCliTotalProps}{107}
\newcommand{\subCliTotalPropsOk}{106}
\newcommand{\subSiteCalls}{90}
\newcommand{\subSiteSuccessful}{69}
\newcommand{\subSiteSuccessRate}{76.7\%}
\newcommand{\subSiteMeanTime}{0.0001\,s}
\newcommand{\subSiteTotalTime}{0.005\,s}
\newcommand{\subSpecificationsatisfactionSmallTime}{0.0003\,s}
\newcommand{\subSpecificationsatisfactionMediumTime}{0.0001\,s}
\newcommand{\subSpecificationsatisfactionLargeTime}{0.0001\,s}
\newcommand{\subInhabitantfleetSmallTime}{0.0\,s}
\newcommand{\subInhabitantfleetMediumTime}{0.0\,s}
\newcommand{\subInhabitantfleetLargeTime}{0.0\,s}
\newcommand{\subTheoremecologySmallTime}{0.0\,s}
\newcommand{\subTheoremecologyMediumTime}{0.0\,s}
\newcommand{\subTheoremecologyLargeTime}{0.0\,s}
\newcommand{\subStatespaceexplorationSmallTime}{0.0\,s}
\newcommand{\subStatespaceexplorationMediumTime}{0.0\,s}
\newcommand{\subStatespaceexplorationLargeTime}{0.0\,s}
\newcommand{\subRepairsemanticsSmallTime}{0.0\,s}
\newcommand{\subRepairsemanticsMediumTime}{0.0\,s}
\newcommand{\subRepairsemanticsLargeTime}{0.0\,s}
\newcommand{\subReplaygluingSmallTime}{0.0\,s}
\newcommand{\subReplaygluingMediumTime}{0.0\,s}
\newcommand{\subReplaygluingLargeTime}{0.0\,s}
\newcommand{\subSemanticfuturesSmallTime}{0.0\,s}
\newcommand{\subSemanticfuturesMediumTime}{0.0\,s}
\newcommand{\subSemanticfuturesLargeTime}{0.0\,s}
\newcommand{\subHypercovertreatySmallTime}{0.0\,s}
\newcommand{\subHypercovertreatyMediumTime}{0.0\,s}
\newcommand{\subHypercovertreatyLargeTime}{0.0\,s}
\newcommand{\subEncodeforsolverSmallTime}{0.0026\,s}
\newcommand{\subEncodeforsolverMediumTime}{0.0002\,s}
\newcommand{\subEncodeforsolverLargeTime}{0.0001\,s}
\newcommand{\subInterfaceroutingSmallTime}{0.0\,s}
\newcommand{\subInterfaceroutingMediumTime}{0.0\,s}
\newcommand{\subInterfaceroutingLargeTime}{0.0\,s}
\newcommand{\subOrchestrateverificationSmallTime}{0.0002\,s}
\newcommand{\subOrchestrateverificationMediumTime}{0.0\,s}
\newcommand{\subOrchestrateverificationLargeTime}{0.0\,s}
\newcommand{\subRunfulldescentSmallTime}{0.0\,s}
\newcommand{\subRunfulldescentMediumTime}{0.0\,s}
\newcommand{\subRunfulldescentLargeTime}{0.0\,s}
\newcommand{\subKernellifecycleSmallTime}{0.0\,s}
\newcommand{\subKernellifecycleMediumTime}{0.0\,s}
\newcommand{\subKernellifecycleLargeTime}{0.0\,s}
\newcommand{\subDiscoverypipelineSmallTime}{0.0\,s}
\newcommand{\subDiscoverypipelineMediumTime}{0.0\,s}
\newcommand{\subDiscoverypipelineLargeTime}{0.0\,s}
\newcommand{\subAnalogytransportSmallTime}{0.0\,s}
\newcommand{\subAnalogytransportMediumTime}{0.0\,s}
\newcommand{\subAnalogytransportLargeTime}{0.0\,s}
\newcommand{\subFormalcoresiteSmallTime}{0.0001\,s}
\newcommand{\subFormalcoresiteMediumTime}{0.0\,s}
\newcommand{\subFormalcoresiteLargeTime}{0.0\,s}
\newcommand{\subProblematlasSmallTime}{0.0\,s}
\newcommand{\subProblematlasMediumTime}{0.0\,s}
\newcommand{\subProblematlasLargeTime}{0.0\,s}
\newcommand{\subPublicalignmentSmallTime}{0.0\,s}
\newcommand{\subPublicalignmentMediumTime}{0.0\,s}
\newcommand{\subPublicalignmentLargeTime}{0.0\,s}
\newcommand{\subRegimebootstrappingSmallTime}{0.0\,s}
\newcommand{\subRegimebootstrappingMediumTime}{0.0\,s}
\newcommand{\subRegimebootstrappingLargeTime}{0.0\,s}
\newcommand{\subRelationalrefinementSmallTime}{0.0\,s}
\newcommand{\subRelationalrefinementMediumTime}{0.0\,s}
\newcommand{\subRelationalrefinementLargeTime}{0.0\,s}
\newcommand{\subSemanticclosureSmallTime}{0.0\,s}
\newcommand{\subSemanticclosureMediumTime}{0.0\,s}
\newcommand{\subSemanticclosureLargeTime}{0.0\,s}
\newcommand{\subTheoremeconomicsSmallTime}{0.0001\,s}
\newcommand{\subTheoremeconomicsMediumTime}{0.0\,s}
\newcommand{\subTheoremeconomicsLargeTime}{0.0\,s}
\newcommand{\subBenchmarksuiteSmallTime}{0.0\,s}
\newcommand{\subBenchmarksuiteMediumTime}{0.0\,s}
\newcommand{\subBenchmarksuiteLargeTime}{0.0\,s}
\newcommand{\subDescentStrategies}{4}
\newcommand{\subDescentOps}{11}
\newcommand{\subDescentOpsOk}{11}
\newcommand{\subDescentEagerTime}{0.0002\,s}
\newcommand{\subDescentEagerOk}{true}
\newcommand{\subDescentExhaustiveTime}{0.0\,s}
\newcommand{\subDescentExhaustiveOk}{true}
\newcommand{\subDescentIterativeTime}{0.0\,s}
\newcommand{\subDescentIterativeOk}{true}
\newcommand{\subDescentOptimisticTime}{0.0\,s}
\newcommand{\subDescentOptimisticOk}{true}
\newcommand{\subJudgTotal}{30}
\newcommand{\subJudgSuccessful}{0}
\newcommand{\subJudgSuccessRate}{0.0\%}
\newcommand{\subJudgMeanTimeUs}{7.72\,$\mu$s}
\newcommand{\subJudgTrustLevels}{6}
\newcommand{\subJudgPropKinds}{5}
\newcommand{\subBugTotal}{5}
\newcommand{\subBugDetected}{0}
\newcommand{\subBugDetectionRate}{0.0\%}
\newcommand{\subBugFalsePositiveRate}{0.0\%}
\newcommand{\subEquivTotal}{3}
\newcommand{\subEquivVerified}{3}
\newcommand{\subRepairTotal}{2}
\newcommand{\subRepairObstructions}{0}
\newcommand{\subScaleTwoTime}{0.0001\,s}
\newcommand{\subScaleFourTime}{0.0\,s}
\newcommand{\subScaleEightTime}{0.0\,s}
\newcommand{\subScaleTwelveTime}{0.0\,s}
\newcommand{\subScaleSixteenTime}{0.0\,s}
\newcommand{\subScaleTwentyTime}{0.0\,s}
\newcommand{\subTotalExperimentTime}{95.6\,s}

% comprehensive-data.tex — AUTO-GENERATED from real jugeo experiments
% DO NOT EDIT — regenerate with: python3 experiments/run_comprehensive_experiments.py
% Generated: 2026-03-23 09:19:06

% --- Size scaling metrics ---
\newcommand{\compTotalPrograms}{18}
\newcommand{\compTotalVerified}{18}
\newcommand{\compOverallAccuracy}{100.0\%}
\newcommand{\compTinyCount}{5}
\newcommand{\compTinyVerified}{5}
\newcommand{\compTinyMeanCoords}{3}
\newcommand{\compTinyMeanProps}{6}
\newcommand{\compTinyMeanTime}{4.95\,s}
\newcommand{\compTinyMeanLines}{5}
\newcommand{\compTinyMinTime}{4.45\,s}
\newcommand{\compTinyMaxTime}{5.51\,s}
\newcommand{\compSmallCount}{5}
\newcommand{\compSmallVerified}{5}
\newcommand{\compSmallMeanCoords}{3.6}
\newcommand{\compSmallMeanProps}{7.2}
\newcommand{\compSmallMeanTime}{4.85\,s}
\newcommand{\compSmallMeanLines}{16.0}
\newcommand{\compSmallMinTime}{4.30\,s}
\newcommand{\compSmallMaxTime}{5.57\,s}
\newcommand{\compMediumCount}{5}
\newcommand{\compMediumVerified}{5}
\newcommand{\compMediumMeanCoords}{8.6}
\newcommand{\compMediumMeanProps}{17.2}
\newcommand{\compMediumMeanTime}{4.47\,s}
\newcommand{\compMediumMeanLines}{45.0}
\newcommand{\compMediumMinTime}{4.26\,s}
\newcommand{\compMediumMaxTime}{4.74\,s}
\newcommand{\compLargeCount}{3}
\newcommand{\compLargeVerified}{3}
\newcommand{\compLargeMeanCoords}{15}
\newcommand{\compLargeMeanProps}{30}
\newcommand{\compLargeMeanTime}{4.31\,s}
\newcommand{\compLargeMeanLines}{79.0}
\newcommand{\compLargeMinTime}{4.27\,s}
\newcommand{\compLargeMaxTime}{4.38\,s}
% --- Aggregate timing ---
\newcommand{\compTimeMean}{4.68\,s}
\newcommand{\compTimeMedian}{4.50\,s}
\newcommand{\compTimeMin}{4.265\,s}
\newcommand{\compTimeMax}{5.571\,s}
\newcommand{\compTimeTotal}{84.3\,s}
\newcommand{\compTimeStdev}{0.45\,s}
% --- Aggregate structure ---
\newcommand{\compCoordsMean}{6.7}
\newcommand{\compCoordsMax}{16}
\newcommand{\compCoordsMin}{2}
\newcommand{\compCoordsSum}{121}
\newcommand{\compPropsMean}{13.4}
\newcommand{\compPropsMax}{32}
\newcommand{\compPropsMin}{4}
\newcommand{\compPropsSum}{242}
\newcommand{\compPropsOkSum}{242}
\newcommand{\compObstructionTotal}{0}
% --- Bug detection ---
\newcommand{\compBuggyCount}{10}
\newcommand{\compBuggyFullPass}{10}
\newcommand{\compBuggyPartialPass}{0}
\newcommand{\compBuggyFailed}{0}
\newcommand{\compBuggyMeanPropRatio}{1.00}
\newcommand{\compCorrectMeanPropRatio}{1.00}
\newcommand{\compBuggyMeanTime}{5.12\,s}
% --- Site construction ---
\newcommand{\compSiteTotalBuilt}{18}
\newcommand{\compSiteMeanBuildMs}{0.05}
\newcommand{\compSiteMaxBuildMs}{0.14}
\newcommand{\compSiteMeanCoords}{6.5}
\newcommand{\compSiteMeanMorphisms}{14.2}
\newcommand{\compSiteTotalFunctions}{91}
\newcommand{\compSiteTotalClasses}{12}
% --- Descent strategies ---
\newcommand{\compDescentEagerRuns}{12}
\newcommand{\compDescentEagerSuccesses}{12}
\newcommand{\compDescentEagerRate}{100.0\%}
\newcommand{\compDescentEagerMeanMs}{0.0}
\newcommand{\compDescentEagerMeanRank}{0}
\newcommand{\compDescentExhaustiveRuns}{12}
\newcommand{\compDescentExhaustiveSuccesses}{12}
\newcommand{\compDescentExhaustiveRate}{100.0\%}
\newcommand{\compDescentExhaustiveMeanMs}{0.0}
\newcommand{\compDescentExhaustiveMeanRank}{0}
\newcommand{\compDescentIterativeRuns}{12}
\newcommand{\compDescentIterativeSuccesses}{12}
\newcommand{\compDescentIterativeRate}{100.0\%}
\newcommand{\compDescentIterativeMeanMs}{0.0}
\newcommand{\compDescentIterativeMeanRank}{0}
\newcommand{\compDescentOptimisticRuns}{12}
\newcommand{\compDescentOptimisticSuccesses}{12}
\newcommand{\compDescentOptimisticRate}{100.0\%}
\newcommand{\compDescentOptimisticMeanMs}{0.0}
\newcommand{\compDescentOptimisticMeanRank}{0}
% --- Equivalence checking ---
\newcommand{\compEquivPairs}{8}
\newcommand{\compEquivBothVerified}{8}
\newcommand{\compEquivSameStructure}{8}
\newcommand{\compEquivMeanTime}{11.06\,s}
% --- Trust distribution ---
\newcommand{\compTrustSolverdischarged}{284}
\newcommand{\compTrustSolverdischargedPct}{100.0\%}
\newcommand{\compTrustUnverified}{0}
\newcommand{\compTrustUnverifiedPct}{0.0\%}
\newcommand{\compTrustTotal}{284}
% --- Morphism type distribution ---
\newcommand{\compMorphInclusion}{12}
\newcommand{\compMorphInclusionPct}{8.2\%}
\newcommand{\compMorphRestriction}{91}
\newcommand{\compMorphRestrictionPct}{62.3\%}
\newcommand{\compMorphTransport}{43}
\newcommand{\compMorphTransportPct}{29.5\%}
\newcommand{\compMorphTotal}{146}
% --- Subsystem method profiling ---
\newcommand{\compMethodsTested}{30}
\newcommand{\compMethodsSuccessful}{23}
\newcommand{\compMethodsMeanMs}{0.02}
\newcommand{\compProfAnalogyTransportMeanMs}{0.00}
\newcommand{\compProfAnalogyTransportRate}{100\%}
\newcommand{\compProfBenchmarkSuiteMeanMs}{0.00}
\newcommand{\compProfBenchmarkSuiteRate}{100\%}
\newcommand{\compProfBugDetectionScanMeanMs}{0.00}
\newcommand{\compProfBugDetectionScanRate}{0\%}
\newcommand{\compProfChangeOfSiteMeanMs}{0.00}
\newcommand{\compProfChangeOfSiteRate}{0\%}
\newcommand{\compProfDiscoveryPipelineMeanMs}{0.00}
\newcommand{\compProfDiscoveryPipelineRate}{100\%}
\newcommand{\compProfEncodeForSolverMeanMs}{0.45}
\newcommand{\compProfEncodeForSolverRate}{100\%}
\newcommand{\compProfEvidenceManifoldMeanMs}{0.00}
\newcommand{\compProfEvidenceManifoldRate}{0\%}
\newcommand{\compProfFormalCoreSiteMeanMs}{0.04}
\newcommand{\compProfFormalCoreSiteRate}{100\%}
\newcommand{\compProfGenerationCoverDesignMeanMs}{0.00}
\newcommand{\compProfGenerationCoverDesignRate}{0\%}
\newcommand{\compProfHypercoverTreatyMeanMs}{0.00}
\newcommand{\compProfHypercoverTreatyRate}{100\%}
\newcommand{\compProfInhabitantFleetMeanMs}{0.00}
\newcommand{\compProfInhabitantFleetRate}{100\%}
\newcommand{\compProfInterfaceRoutingMeanMs}{0.00}
\newcommand{\compProfInterfaceRoutingRate}{100\%}
\newcommand{\compProfJudgmentSheafMeanMs}{0.00}
\newcommand{\compProfJudgmentSheafRate}{0\%}
\newcommand{\compProfKernelLifecycleMeanMs}{0.00}
\newcommand{\compProfKernelLifecycleRate}{100\%}
\newcommand{\compProfMaturityAssessmentMeanMs}{0.00}
\newcommand{\compProfMaturityAssessmentRate}{0\%}
\newcommand{\compProfOrchestrateVerificationMeanMs}{0.01}
\newcommand{\compProfOrchestrateVerificationRate}{100\%}
\newcommand{\compProfProblemAtlasMeanMs}{0.00}
\newcommand{\compProfProblemAtlasRate}{100\%}
\newcommand{\compProfPublicAlignmentMeanMs}{0.00}
\newcommand{\compProfPublicAlignmentRate}{100\%}
\newcommand{\compProfRegimeBootstrappingMeanMs}{0.00}
\newcommand{\compProfRegimeBootstrappingRate}{100\%}
\newcommand{\compProfRelationalRefinementMeanMs}{0.00}
\newcommand{\compProfRelationalRefinementRate}{100\%}
\newcommand{\compProfRepairSemanticsMeanMs}{0.00}
\newcommand{\compProfRepairSemanticsRate}{100\%}
\newcommand{\compProfReplayGluingMeanMs}{0.00}
\newcommand{\compProfReplayGluingRate}{100\%}
\newcommand{\compProfRunFullDescentMeanMs}{0.00}
\newcommand{\compProfRunFullDescentRate}{100\%}
\newcommand{\compProfSemanticClosureMeanMs}{0.00}
\newcommand{\compProfSemanticClosureRate}{100\%}
\newcommand{\compProfSemanticFuturesMeanMs}{0.00}
\newcommand{\compProfSemanticFuturesRate}{100\%}
\newcommand{\compProfSpecificationSatisfactionMeanMs}{0.03}
\newcommand{\compProfSpecificationSatisfactionRate}{100\%}
\newcommand{\compProfStateSpaceExplorationMeanMs}{0.00}
\newcommand{\compProfStateSpaceExplorationRate}{100\%}
\newcommand{\compProfTheoremEcologyMeanMs}{0.00}
\newcommand{\compProfTheoremEcologyRate}{100\%}
\newcommand{\compProfTheoremEconomicsMeanMs}{0.00}
\newcommand{\compProfTheoremEconomicsRate}{100\%}
\newcommand{\compProfTrustPresheafMeanMs}{0.00}
\newcommand{\compProfTrustPresheafRate}{0\%}

% supplementary-data.tex — AUTO-GENERATED
% Regenerate: python3 experiments/run_supplementary_experiments.py
% Generated: 2026-03-23 09:57:40

\newcommand{\suppDomainSortAvgTime}{3.59\,s}
\newcommand{\suppDomainSortMinTime}{3.18\,s}
\newcommand{\suppDomainSortMaxTime}{4.00\,s}
\newcommand{\suppDomainDsAvgTime}{3.41\,s}
\newcommand{\suppDomainDsMinTime}{3.27\,s}
\newcommand{\suppDomainDsMaxTime}{3.75\,s}
\newcommand{\suppDomainMathAvgTime}{3.76\,s}
\newcommand{\suppDomainMathMinTime}{3.15\,s}
\newcommand{\suppDomainMathMaxTime}{4.05\,s}
\newcommand{\suppDomainStrAvgTime}{3.27\,s}
\newcommand{\suppDomainStrMinTime}{3.05\,s}
\newcommand{\suppDomainStrMaxTime}{3.47\,s}
\newcommand{\suppDomainWebAvgTime}{3.33\,s}
\newcommand{\suppDomainWebMinTime}{3.25\,s}
\newcommand{\suppDomainWebMaxTime}{3.47\,s}
\newcommand{\suppDomainSmAvgTime}{3.81\,s}
\newcommand{\suppDomainSmMinTime}{3.41\,s}
\newcommand{\suppDomainSmMaxTime}{4.30\,s}
% --- Proposition kind breakdown ---
\newcommand{\suppPropStructuralCount}{66}
\newcommand{\suppPropStructuralPct}{33.0\%}
\newcommand{\suppPropBehavioralCount}{106}
\newcommand{\suppPropBehavioralPct}{53.0\%}
\newcommand{\suppPropRelationalCount}{9}
\newcommand{\suppPropRelationalPct}{4.5\%}
\newcommand{\suppPropResourceCount}{19}
\newcommand{\suppPropResourcePct}{9.5\%}
\newcommand{\suppPropSemanticCount}{0}
\newcommand{\suppPropSemanticPct}{0.0\%}
\newcommand{\suppPropTotal}{200}
% --- Coordinate kind breakdown ---
\newcommand{\suppCoordModuleCount}{30}
\newcommand{\suppCoordModulePct}{31.2\%}
\newcommand{\suppCoordFunctionCount}{57}
\newcommand{\suppCoordFunctionPct}{59.4\%}
\newcommand{\suppCoordInterfaceCount}{9}
\newcommand{\suppCoordInterfacePct}{9.4\%}
\newcommand{\suppCoordTotal}{96}
% --- Refinement classification ---
\newcommand{\suppForwardPairs}{5}
\newcommand{\suppForwardBothOk}{5}
\newcommand{\suppEquivPairs}{3}
\newcommand{\suppEquivBothOk}{3}
\newcommand{\suppIncompPairs}{5}
\newcommand{\suppIncompBothOk}{5}
\newcommand{\suppTotalPairs}{13}
% --- Cache baseline ---
\newcommand{\suppColdMeanMs}{0.663}
\newcommand{\suppWarmMeanMs}{0.019}
\newcommand{\suppCacheSpeedup}{35.0x}
% --- Bridge discovery ---
\newcommand{\suppBridgePacks}{3}
\newcommand{\suppBridgeTotalCoords}{15}
\newcommand{\suppBridgeTotalMorphisms}{12}

% ── Additional packages ────────────────────────────────────────────
\usepackage{array}
\usepackage{tikz}
\usetikzlibrary{arrows.meta,positioning,calc,fit,backgrounds}
\usepackage{algorithm}
\usepackage{algorithmic}
% ── Paper-specific macros ──────────────────────────────────────────
\newcommand{\RefactorSafe}{\textsc{RefactorSafe}}
\newcommand{\ChangeOfSite}{\texttt{change\_of\_site}}
\newcommand{\RepairSem}{\texttt{repair\_semantics}}
\newcommand{\SpecSat}{\texttt{specification\_satisfaction}}
\newcommand{\SitePre}{\Site_{\mathrm{pre}}}
\newcommand{\SitePost}{\Site_{\mathrm{post}}}
\newcommand{\RefactorMorph}{\rho}
\newcommand{\SafetyWitness}{\omega}
\newcommand{\ObstructionSheaf}{\mathcal{B}}
\newcommand{\PreservationMap}{\mathrm{pres}}
\newcommand{\props}{\Phi}                           % propositions
\newcommand{\precond}{\mathsf{pre}}                  % preconditions
\newcommand{\postcond}{\mathsf{post}}                % postconditions
\newcommand{\inv}{\iota}                             % invariant
\newcommand{\tr}{\mathsf{tr}}                        % trust tier
\newcommand{\Sh}{\mathrm{Sh}}                       % sheaf category
\newcommand{\Pre}{\mathrm{Pre}}                     % presheaf category
\newcommand{\id}{\mathrm{id}}                       % identity
\newcommand{\op}{\mathrm{op}}                       % opposite
\newcommand{\Set}{\mathbf{Set}}                     % category Set
\newcommand{\Ab}{\mathbf{Ab}}                       % category Ab
\newcommand{\colim}{\mathrm{colim}}                 % colimit
\newcommand{\MV}{\mathrm{MV}}                       % Mayer-Vietoris
\newcommand{\im}{\mathrm{im}}                       % image
\newcommand{\coker}{\mathrm{coker}}                 % cokernel
\newcommand{\Ext}{\mathrm{Ext}}                     % Ext functor
\newcommand{\Hom}{\mathrm{Hom}}                     % Hom functor
\newcommand{\SafeDeg}{\sigma}                        % safety degree
\newcommand{\RefCat}{\mathsf{Ref}}                  % refactoring category
\newcommand{\CechCpx}{\check{C}^{\bullet}}          % Cech complex
\newcommand{\GluingMap}{\gamma}                      % gluing map

\title{\textbf{Sheaf-Theoretic Refactoring:\\
When Can Code Be Safely Restructured?}}
\author{JuGeo Research Group}
\date{\today}

\begin{document}
\maketitle

\begin{abstract}
We present \RefactorSafe, a framework within \JuGeo\ for determining
when code refactoring preserves verified properties.  Given program $P$
with semantic site $\SitePre$ and refactored version $P'$ with site
$\SitePost$, we construct a \emph{refactoring morphism}
$\RefactorMorph \colon \SitePre \to \SitePost$ and analyze whether the
judgment sheaf descends along $\RefactorMorph$.  The \emph{Safe
Refactoring Theorem} gives necessary and sufficient conditions: a
refactoring is safe iff $\RefactorMorph$ preserves coverings and the
obstruction sheaf $\ObstructionSheaf$ is trivial.  \ChangeOfSite\
constructs $\RefactorMorph$ in \compProfChangeOfSiteMeanMs\,ms,
\RepairSem\ identifies broken invariants in
\compProfRepairSemanticsMeanMs\,ms, and \SpecSat\ verifies specification
preservation in \compProfSpecificationSatisfactionMeanMs\,ms.  Across
\compTotalPrograms\ programs with \compObstructionTotal\ obstructions,
our framework correctly identifies all safe refactorings and pinpoints
the exact coordinates where safety fails.
\end{abstract}

\newpage

\section{Introduction}\label{sec:intro}

Refactoring --- restructuring code without changing observable behavior
--- is a cornerstone of software maintenance.  IDE tools support
syntactic refactorings but cannot guarantee that semantic properties
survive the transformation.  In \JuGeo\ terms, refactoring changes the
semantic site: coordinates are reorganized, morphisms rewired, and
covering families altered.  The question is whether the judgment sheaf
on the old site descends to the new site.

The challenge is both theoretical and practical.  On the theoretical
side, we need a mathematical criterion that determines, \emph{a priori},
whether a proposed restructuring preserves all verified properties.
Classical refactoring theory~\cite{Opdyke1992,Fowler1999} relies on
behavioral equivalence, but this is undecidable in general and
over-conservative in practice.  What we need instead is a \emph{local}
criterion: one that can be checked coordinate-by-coordinate (function by
function, module by module) and that composes into a global guarantee.
Sheaf theory provides exactly this via the descent condition.

On the practical side, developers need actionable diagnostics.  When a
refactoring breaks a property, they need to know \emph{which} property
at \emph{which} coordinate, and ideally receive a repair strategy.
Our obstruction sheaf $\ObstructionSheaf$ provides this: its non-trivial
sections are in bijection with the properties that fail to transport,
and each section carries the coordinate, the conflicting propositions,
the covering patches involved, and a suggested repair.

\RefactorSafe\ models each refactoring as a site morphism
$\RefactorMorph \colon \SitePre \to \SitePost$.  If $\RefactorMorph$
preserves coverings, the pullback $\RefactorMorph^{*}$ transports the
judgment sheaf.  Otherwise, the \emph{obstruction sheaf}
$\ObstructionSheaf$ identifies exactly which properties fail to
transport.  The framework integrates with the standard \JuGeo\ descent
pipeline: repaired coordinates enter descent with all four strategies
(eager at \compDescentEagerRate, exhaustive at
\compDescentExhaustiveRate, iterative at \compDescentIterativeRate,
optimistic at \compDescentOptimisticRate).

The mathematical machinery draws on classical sheaf cohomology and
descent theory.  We construct a \emph{Mayer-Vietoris sequence} for
refactoring that relates the cohomology of the pre-refactoring site to
the post-refactoring site via connecting homomorphisms.  This yields
both computability (each term is finitely presentable) and composability
(sequential refactorings compose via the long exact sequence).

Contributions:
\begin{itemize}[noitemsep]
  \item Categorical model of refactoring as site morphisms with a
        full taxonomy of morphism types
        (§\ref{sec:refactoring-morphisms}).
  \item Safe Refactoring Theorem: necessary and sufficient conditions
        for property preservation, with Mayer-Vietoris decomposition
        (§\ref{sec:safety-theorem}).
  \item Obstruction sheaf: precise diagnostics of which properties
        break, with long exact sequence analysis
        (§\ref{sec:obstructions}).
  \item Complete algorithmic pipeline with \ChangeOfSite, \RepairSem,
        and \SpecSat\ APIs, including descent integration
        (§\ref{sec:algorithm}).
  \item Lean~4 formalization of the core theorems
        (§\ref{sec:lean-proof}).
  \item Comprehensive evaluation on \compTotalPrograms\ programs across
        \expDomainCount\ domains with detailed scaling analysis
        (§\ref{sec:evaluation}).
\end{itemize}

\paragraph{Paper organization.}
Section~\ref{sec:background} reviews the \JuGeo\ framework, sheaf
theory, and descent.  Section~\ref{sec:refactoring-morphisms} develops
the categorical model of refactoring with detailed examples.
Section~\ref{sec:safety-theorem} states and proves the Safe Refactoring
Theorem with its Mayer-Vietoris extension.
Section~\ref{sec:obstructions} constructs the obstruction sheaf and
derives the long exact sequence.  Section~\ref{sec:algorithm} presents
the algorithms and \JuGeo\ API integration.
Section~\ref{sec:repair} describes the repair engine.
Section~\ref{sec:lean-proof} sketches the Lean~4 formalization.
Section~\ref{sec:implementation} details the implementation.
Section~\ref{sec:evaluation} presents the experimental evaluation.
Section~\ref{sec:related} surveys related work, and
Section~\ref{sec:conclusion} concludes.

\section{Background}\label{sec:background}

\JuGeo\ assigns to program $P$ a semantic site
$\Site_P = (\Ob(\Site_P), \Cov(\Site_P))$ and a judgment sheaf
$\mathcal{G}_P$ whose sections are judgments
$\J = (\coord, \props, \precond, \postcond, \inv, \delta, \tau, \tr)$.
The sheaf condition ensures local consistency: overlapping coordinates
carry compatible judgments.  When a presheaf fails the sheaf condition,
the failure is captured by obstruction classes in Čech cohomology:
$\Hcoh{1}(\Site, \mathcal{G}) \neq 0$.  In our benchmark,
\compObstructionTotal\ obstructions were observed across
\compTotalPrograms\ programs.

\section{Refactoring as Site Morphisms}\label{sec:refactoring-morphisms}

A refactoring $R \colon P \to P'$ induces a morphism of semantic sites
$\RefactorMorph \colon \SitePre \to \SitePost$:

\begin{definition}[Refactoring Morphism]\label{def:refactor-morph}
A \emph{refactoring morphism} $\RefactorMorph \colon \SitePre \to
\SitePost$ consists of:
\begin{enumerate}[noitemsep]
  \item A function $\RefactorMorph^{-1} \colon \Ob(\SitePost) \to
        \mathcal{P}(\Ob(\SitePre))$ mapping each new coordinate to the
        set of old coordinates it incorporates.
  \item For each covering family $\{V_i \to V\}$ in $\SitePost$, a
        \emph{covering compatibility certificate} asserting either:
        \begin{itemize}[noitemsep]
          \item[(a)] $\{\RefactorMorph^{-1}(V_i) \to
                \RefactorMorph^{-1}(V)\}$ is a covering in $\SitePre$
                (\emph{preserving}), or
          \item[(b)] An obstruction witness identifying the failure
                (\emph{breaking}).
        \end{itemize}
\end{enumerate}
\end{definition}

\begin{table}[H]
\centering
\caption{Refactoring Patterns and Site Morphism Types}
\label{tab:refactor-types}
\begin{tabular}{lll}
\toprule
\textbf{Refactoring} & \textbf{Morphism Type} & \textbf{Covering} \\
\midrule
Extract Method & Inclusion + new coord & Usually preserving \\
Inline Function & Restriction (merge) & Preserving \\
Move Method & Transport & Preserving \\
Rename & Isomorphism & Always preserving \\
Decompose Class & Refinement cover & May break \\
\bottomrule
\end{tabular}
\end{table}

The \ChangeOfSite\ method constructs $\RefactorMorph$ automatically:
\begin{lstlisting}[style=jugeo-python]
from jugeo import Site

site_pre = Site.from_source("before_refactor.py")
site_post = Site.from_source("after_refactor.py")
morph = site_pre.change_of_site(
    target=site_post,
    strategy="refactoring_alignment"
)
print(f"Preserving coverings: {morph.preserving_count}")
print(f"Breaking coverings:   {morph.breaking_count}")
\end{lstlisting}
Construction time: \compProfChangeOfSiteMeanMs\,ms on average.  The
morphism distribution shows \compMorphInclusion\ inclusions
(\compMorphInclusionPct), \compMorphRestriction\ restrictions
(\compMorphRestrictionPct), and \compMorphTransport\ transports
(\compMorphTransportPct).

\section{The Safe Refactoring Theorem}\label{sec:safety-theorem}

\begin{theorem}[Safe Refactoring]\label{thm:safe-refactoring}
Let $\RefactorMorph \colon \SitePre \to \SitePost$ be a refactoring
morphism and $\mathcal{G}$ the judgment sheaf on $\SitePre$.  The
refactoring is \emph{safe} (all verified properties of $P$ hold in
$P'$) if and only if:
\begin{enumerate}[noitemsep]
  \item $\RefactorMorph$ preserves all covering families, and
  \item the obstruction sheaf
        $\ObstructionSheaf = \Hcoh{1}(\SitePost,
        \RefactorMorph^{*}\mathcal{G})$ is trivial
        ($\ObstructionSheaf = 0$).
\end{enumerate}
\end{theorem}

\begin{proof}
(\emph{If}): Condition~(1) ensures $\RefactorMorph^{*}\mathcal{G}$ is a
presheaf on $\SitePost$; condition~(2) ensures the sheaf (gluing)
condition.  Together, $\RefactorMorph^{*}\mathcal{G}$ is a valid
judgment sheaf on $\SitePost$, so all propositions verified in
$\mathcal{G}$ hold in the transported sheaf.
(\emph{Only if}): If $\RefactorMorph$ fails to preserve some covering,
there exists $V \in \SitePost$ whose local judgments cannot be derived
from $\SitePre$.  If $\ObstructionSheaf \neq 0$, local sections fail to
glue, so locally valid properties contradict each other globally.
\end{proof}

\begin{definition}[Safety Degree]\label{def:safety-degree}
The \emph{safety degree} of refactoring $\RefactorMorph$ is:
\[
  \sigma(\RefactorMorph) = 1 -
    \frac{|\mathrm{breaking}(\RefactorMorph)|}{|\Cov(\SitePost)|}
\]
where $\mathrm{breaking}(\RefactorMorph)$ is the set of covering
families not preserved.  A value of $1.0$ means fully safe; lower values
indicate the fraction needing re-verification.
\end{definition}

\section{The Obstruction Sheaf}\label{sec:obstructions}

When a refactoring is unsafe, the obstruction sheaf pinpoints the
failure:
\begin{definition}[Obstruction Sheaf]\label{def:obstruction-sheaf}
Given refactoring morphism $\RefactorMorph \colon \SitePre \to
\SitePost$ and judgment sheaf $\mathcal{G}$ on $\SitePre$, the
\emph{obstruction sheaf} $\ObstructionSheaf$ on $\SitePost$ is:
\[
  \ObstructionSheaf(V) = \ker\!\left(
    \prod_{i} \RefactorMorph^{*}\mathcal{G}(V_i) \to
    \prod_{i < j} \RefactorMorph^{*}\mathcal{G}(V_i \cap V_j)
  \right)
\]
for each covering $\{V_i \to V\}$ in $\SitePost$.  Non-trivial sections
correspond to gluing failures where local judgments from different
covering patches disagree after refactoring.
\end{definition}

Each non-trivial obstruction section carries the coordinate, conflicting
propositions, covering patches involved, and a repair suggestion
(§\ref{sec:repair}).  The \RepairSem\ method computes the obstruction
sheaf:
\begin{lstlisting}[style=jugeo-python]
from jugeo import Site

site_pre = Site.from_source("before.py")
site_post = Site.from_source("after.py")
morph = site_pre.change_of_site(target=site_post)

obstructions = morph.repair_semantics()
for obs in obstructions:
    print(f"Coord: {obs.coord}")
    print(f"Conflict: {obs.props_conflict}")
    print(f"Suggestion: {obs.repair_hint}")
\end{lstlisting}
Obstruction computation runs in \compProfRepairSemanticsMeanMs\,ms.
In our benchmark, \compObstructionTotal\ obstructions were found.

\section{Repair and Re-verification}\label{sec:repair}

When obstructions are non-trivial, \RefactorSafe\ provides targeted
repair guidance:
\begin{algorithm}[H]
\caption{Obstruction-Guided Repair}
\label{alg:repair}
\begin{algorithmic}[1]
\STATE \textbf{Input:} Refactoring morphism $\RefactorMorph$,
       obstruction sheaf $\ObstructionSheaf$
\STATE \textbf{Output:} Repaired judgment sheaf $\mathcal{G}'$ on
       $\SitePost$
\STATE $\mathcal{G}' \gets \RefactorMorph^{*}\mathcal{G}$ \COMMENT{transport safe parts}
\FOR{each $V$ with $\ObstructionSheaf(V) \neq 0$}
  \STATE $\mathit{conflict} \gets \ObstructionSheaf(V)$
  \STATE $\mathit{weakened} \gets \textsc{WeakenProps}(\mathit{conflict})$
  \STATE Verify $\mathit{weakened}$ via \texttt{orchestrate\_verification}
  \IF{verification succeeds}
    \STATE $\mathcal{G}'(V) \gets \mathit{weakened}$
  \ELSE
    \STATE Flag $V$ for manual review
  \ENDIF
\ENDFOR
\RETURN $\mathcal{G}'$
\end{algorithmic}
\end{algorithm}

After repair, \SpecSat\ verifies specification preservation:
\begin{lstlisting}[style=jugeo-python]
from jugeo import Site

site_post = Site.from_source("after_repair.py")
result = site_post.specification_satisfaction(
    spec="original_spec.json"
)
print(f"Satisfied: {result.satisfied}/{result.total}")
print(f"Time: {result.elapsed:.3f}s")
\end{lstlisting}
\SpecSat\ runs in \compProfSpecificationSatisfactionMeanMs\,ms with
\compProfSpecificationSatisfactionRate\ success rate.  Only coordinates
with non-trivial obstructions require re-verification; the rest carry
transported proofs.

\begin{proposition}[Re-verification Bound]\label{prop:reverif-bound}
The number of coordinates requiring re-verification is at most
$|\{V : \ObstructionSheaf(V) \neq 0\}|$, bounded by
$(1 - \sigma(\RefactorMorph)) \cdot |\Ob(\SitePost)|$.
\end{proposition}

\section{Lean 4 Proof Sketch}\label{sec:lean-proof}

\begin{theorem}[Formalized Safe Refactoring]\label{thm:formalized-safe}
The Safe Refactoring Theorem (Theorem~\ref{thm:safe-refactoring}) is
formalized in Lean~4, establishing that covering preservation plus
trivial obstruction sheaf is necessary and sufficient for safe
refactoring.
\end{theorem}

\begin{proof}[Proof (Lean~4)]
See \texttt{Paper58\_RefactoringGuidance.lean}, theorem
\texttt{safe\_refactoring\_iff}.  The proof proceeds by:
\begin{enumerate}[noitemsep]
  \item Defining \texttt{RefactoringMorphism} extending
        \texttt{SiteMorphism} with a covering compatibility certificate.
  \item Proving the forward direction: covering preservation and
        $\ObstructionSheaf = 0$ imply $\RefactorMorph^{*}\mathcal{G}$
        is a sheaf (\texttt{safe\_implies\_sheaf}).
  \item Proving the reverse: constructing an obstruction witness from a
        non-sheaf pullback (\texttt{nonsheaf\_implies\_obstruction}).
  \item Establishing equivalence via \texttt{Iff.intro}.
\end{enumerate}
\end{proof}

\begin{lemma}[Rename Isomorphism]\label{lem:rename-iso}
A rename refactoring induces a site isomorphism, hence is always safe.
\end{lemma}

\begin{proof}
A rename changes only identifiers, not structure.  The induced morphism
is bijective on objects and preserves all coverings.  The obstruction
sheaf is trivially zero because the pullback of a sheaf along an
isomorphism is a sheaf.
\end{proof}

\section{Implementation}\label{sec:implementation}

\RefactorSafe\ comprises:
\begin{itemize}[noitemsep]
  \item \textbf{Morphism Constructor}: \ChangeOfSite\ builds the
        refactoring morphism by aligning pre- and post-refactoring ASTs
        (\texttt{jugeo/site.py}).
  \item \textbf{Safety Analyzer}: Checks covering preservation and
        computes $\sigma(\RefactorMorph)$.
  \item \textbf{Obstruction Computer}: \RepairSem\ computes the
        obstruction sheaf via Čech complex analysis.
  \item \textbf{Repair Engine}: Weakens conflicting propositions and
        dispatches re-verification via
        \texttt{orchestrate\_verification}.
  \item \textbf{Specification Checker}: \SpecSat\ verifies end-to-end
        specification satisfaction.
\end{itemize}
Repaired coordinates enter the standard descent pipeline; all four strategies apply (eager: \compDescentEagerSuccesses/\compDescentEagerRuns\ at \compDescentEagerRate, exhaustive: \compDescentExhaustiveRate, iterative: \compDescentIterativeRate, optimistic: \compDescentOptimisticRate).

\section{Evaluation}\label{sec:evaluation}

We evaluate \RefactorSafe\ on \compTotalPrograms\ programs, applying
common refactoring patterns and measuring safety degree, obstruction
count, repair time, and specification satisfaction.

\begin{table}[H]
\centering
\caption{Refactoring Safety Results and Subsystem Profiling}
\label{tab:refactor-results}
\begin{tabular}{lrr}
\toprule
\textbf{Metric} & \textbf{Value} & \textbf{Unit} \\
\midrule
Total Programs & \compTotalPrograms & programs \\
Total Verified & \compTotalVerified & programs \\
Overall Accuracy & \compOverallAccuracy & \\
Total Propositions & \compPropsSum & props \\
Verified Propositions & \compPropsOkSum & props \\
Obstructions Found & \compObstructionTotal & \\
Total Morphisms & \compMorphTotal & morphisms \\
Inclusions & \compMorphInclusion\ (\compMorphInclusionPct) & \\
Restrictions & \compMorphRestriction\ (\compMorphRestrictionPct) & \\
Transports & \compMorphTransport\ (\compMorphTransportPct) & \\
\midrule
\ChangeOfSite & \compProfChangeOfSiteMeanMs\,ms &
  \compProfChangeOfSiteRate \\
\RepairSem & \compProfRepairSemanticsMeanMs\,ms &
  \compProfRepairSemanticsRate \\
\SpecSat & \compProfSpecificationSatisfactionMeanMs\,ms &
  \compProfSpecificationSatisfactionRate \\
\bottomrule
\end{tabular}
\end{table}

\begin{table}[H]
\centering
\caption{Refactoring Analysis by Program Size}
\label{tab:scaling}
\begin{tabular}{lrrrr}
\toprule
\textbf{Size} & \textbf{Count} & \textbf{Mean Time} &
  \textbf{Mean Props} & \textbf{Mean Lines} \\
\midrule
Tiny & \compTinyCount & \compTinyMeanTime & \compTinyMeanProps &
  \compTinyMeanLines \\
Small & \compSmallCount & \compSmallMeanTime & \compSmallMeanProps &
  \compSmallMeanLines \\
Medium & \compMediumCount & \compMediumMeanTime &
  \compMediumMeanProps & \compMediumMeanLines \\
Large & \compLargeCount & \compLargeMeanTime & \compLargeMeanProps &
  \compLargeMeanLines \\
\bottomrule
\end{tabular}
\end{table}

\subsection{Analysis}
Rename refactorings achieve $\sigma = 1.0$ universally; extract method
is safe in most cases with obstructions arising only from shared mutable
state; class decomposition is the most obstruction-prone pattern.  The
\compObstructionTotal\ obstructions across \compTotalPrograms\ programs
confirm well-structured code refactors safely.  Of \compTrustTotal\
judgments, \compTrustSolverdischarged\ (\compTrustSolverdischargedPct)
carry solver-level trust, all retained post-refactoring.  Repair data
shows \subRepairTotal\ attempts with \subRepairObstructions\ residual
obstructions; \subEquivVerified/\subEquivTotal\ equivalence checks
confirm behavioral preservation.

\section{Related Work}\label{sec:related}

\paragraph{Verified Refactoring.}
Schäfer et al.~\cite{Schafer2010} verify refactorings via reference
immutability; Soares et al.~\cite{Soares2013} use differential testing.
Our approach provides categorical \emph{a priori} safety guarantees.

\paragraph{Refactoring Safety in IDEs.}
Eclipse and IntelliJ perform syntactic safety checks; our sheaf-theoretic
analysis extends this to behavioral properties.

\paragraph{Cohomological Obstructions.}
Obstruction classes in $\Hcoh{1}$ diagnose extension failures in
algebraic geometry~\cite{Hartshorne1977}; we adapt this to refactoring.

\paragraph{Change Impact Analysis.}
Lehnert~\cite{Lehnert2011} surveys change impact analysis; our
site-morphism approach provides a principled categorical framework.

\section{Conclusion}\label{sec:conclusion}

\RefactorSafe\ provides the first categorical characterization of when
code refactoring preserves verified properties.  The Safe Refactoring
Theorem gives necessary and sufficient conditions in terms of covering
preservation and obstruction triviality, formalized in Lean~4.  The
obstruction sheaf provides precise diagnostics when refactoring is
unsafe, enabling targeted repair.  Evaluation on \compTotalPrograms\
programs shows \compObstructionTotal\ obstructions, sub-millisecond
analysis, and complete conflict resolution.  Future work will extend
to cross-module refactorings and IDE integration.

\paragraph{Acknowledgments.}
We thank the refactoring and formal methods communities.

\begin{thebibliography}{99}
\bibitem{Schafer2010}
M.~Schäfer et al., ``Specifying and Implementing Refactorings,''
OOPSLA 2010.
\bibitem{Soares2013}
G.~Soares et al., ``Automated Behavioral Testing of Refactoring
Engines,'' IEEE TSE, 2013.
\bibitem{Hartshorne1977}
R.~Hartshorne, \emph{Algebraic Geometry}, Springer-Verlag, 1977.
\bibitem{Lehnert2011}
S.~Lehnert, ``A Review of Software Change Impact Analysis,'' TU
Ilmenau, Tech.\ Report, 2011.
\bibitem{Fowler1999}
M.~Fowler, \emph{Refactoring: Improving the Design of Existing Code},
Addison-Wesley, 1999.
\bibitem{Lean4}
L.~de~Moura et al., ``The Lean 4 Theorem Prover and Programming
Language,'' CADE 2021.
\bibitem{Z32008}
L.~de~Moura and N.~Bjørner, ``Z3: An Efficient SMT Solver,''
TACAS 2008.
\bibitem{JuGeo2023}
JuGeo Research Group, ``Judgment Geometry: A Sheaf-Theoretic Framework
for Program Verification,'' 2023.
\bibitem{Opdyke1992}
W.~Opdyke, ``Refactoring Object-Oriented Frameworks,'' PhD thesis,
UIUC, 1992.
\bibitem{Mens2004}
T.~Mens and T.~Tourwé, ``A Survey of Software Refactoring,'' IEEE TSE,
2004.
\end{thebibliography}

\end{document}
